% Lab 4 of Intelligent Robotics (ASY2): ICP after Besl and McKay on laser scans, alone and with odometry.
% Build: make  (latexmk -pdf). Figures are TikZ files in figures/, made by tools/make_figures.py.
\documentclass[11pt,a4paper,parskip=half]{scrartcl}
\usepackage[english]{babel}
\usepackage[T1]{fontenc}
\usepackage[utf8]{inputenc}
\usepackage[sfdefault,lining]{FiraSans}
\usepackage[scale=0.88]{FiraMono}
\usepackage[a4paper,left=2.4cm,right=2.4cm,top=2.6cm,bottom=2.6cm,headsep=0.8cm]{geometry}
\usepackage{amsmath,amssymb}
\usepackage{booktabs,array,tabularx}
\usepackage{xcolor}
\usepackage{listings}
\usepackage{tikz}
\usetikzlibrary{arrows.meta,calc,shapes.geometric}
\usepackage[skins,breakable]{tcolorbox}
\usepackage{fontawesome5}
\usepackage[font=small,labelfont={bf,color=ohmred},skip=4pt]{caption}
\usepackage{enumitem}
\usepackage{microtype}
\usepackage[headsepline=0.4pt]{scrlayer-scrpage}
\usepackage[colorlinks=true,linkcolor=ohmred,urlcolor=ohmred]{hyperref}
\usepackage{cleveref}

% ---------------------------------------------------------------- colours and headings
\definecolor{ohmred}{HTML}{C8372D}
\definecolor{ohmgreen}{HTML}{2E8B57}
\definecolor{ohmamber}{HTML}{D98E04}
\definecolor{ohmgrey}{HTML}{5A5A5A}
\definecolor{odomblue}{HTML}{4682DC}
\addtokomafont{disposition}{\color{ohmred}}
\addtokomafont{section}{\Large}
\addtokomafont{subsection}{\large}
\setkomafont{pageheadfoot}{\small\color{ohmgrey}}
\setkomafont{pagenumber}{\small\bfseries\color{ohmred}}
\clearpairofpagestyles
\ihead{Intelligent Robotics (ASY2) · Lab 4}
\ohead{ICP}
\ofoot{\pagemark}
\pagestyle{scrheadings}

\lstset{basicstyle=\ttfamily\small, columns=fullflexible, keepspaces=true,
        showstringspaces=false, breaklines=true, frame=none, backgroundcolor=\color{black!5},
        xleftmargin=8pt, xrightmargin=8pt, framexleftmargin=6pt, framexrightmargin=6pt,
        aboveskip=4pt, belowskip=4pt}

\newcommand{\thema}[1]{\texttt{#1}}
\newcommand{\datei}[1]{\texttt{#1}}
\newcommand{\minutes}[1]{\hfill{\normalsize\mdseries\color{ohmgrey}\faClock[regular]~#1 min}}
\newcommand{\matsym}[1]{\mathbf{#1}}
\renewcommand{\vec}[1]{\mathbf{#1}}
\newcommand{\T}{^{\mathsf{T}}}

% ---------------------------------------------------------------- boxes
\tcbset{base/.style={enhanced, breakable, boxrule=0pt, arc=1.5pt, left=8pt, right=8pt,
        top=5pt, bottom=5pt, before skip=8pt, after skip=8pt}}
% \begin{todo}{1.3}: the block STUDENT 1.3, with its hint command at the end of the last line
\newcommand{\hintcmd}[1]{\unskip\nobreak\hfil\penalty50\hskip1em\null\nobreak\hfil
        \mbox{\small\color{ohmgrey}\faLightbulb[regular]~\texttt{python3 hint.py #1}}{\parfillskip=0pt\par}}
\newtcolorbox{todo}[1]{base, colback=ohmred!5, borderline west={2.5pt}{0pt}{ohmred},
        title={\faPen~STUDENT #1}, coltitle=ohmred, colbacktitle=ohmred!5, fonttitle=\bfseries,
        attach title to upper={\quad}, after upper={\hintcmd{#1}}}
\newtcolorbox{script}[1]{base, colback=black!3, borderline west={2.5pt}{0pt}{ohmgrey},
        title={\faBook~Lecture script, #1}, coltitle=ohmgrey, colbacktitle=black!3,
        fonttitle=\small\bfseries, toptitle=2pt, bottomtitle=0pt}
\newtcolorbox{note}{base, colback=ohmamber!8, borderline west={2.5pt}{0pt}{ohmamber}}
\newtcolorbox{checkpoint}{base, colback=ohmgreen!7, borderline west={2.5pt}{0pt}{ohmgreen},
        title={\faCheckCircle[regular]~Checkpoint}, coltitle=ohmgreen, colbacktitle=ohmgreen!7,
        fonttitle=\bfseries, attach title to upper={\quad}}

% \running{<launch arguments>}{<node arguments>}: what terminals 1 and 2 run during a task
\newcommand{\running}[2]{\begin{tcolorbox}[base, colback=black!3, borderline west={2.5pt}{0pt}{ohmgrey},
        title={\faTerminal~Running}, coltitle=ohmgrey, colbacktitle=black!3, fonttitle=\small\bfseries,
        toptitle=2pt, bottomtitle=0pt, fontupper=\small]
  \begin{tabularx}{\linewidth}{@{}>{\color{ohmgrey}}l >{\raggedright\arraybackslash}X@{}}
    terminal 1 & \texttt{ros2 launch ohm\_icp icp.launch.py#1} \\
    terminal 2 & \texttt{ros2 run ohm\_icp icp\_node.py -\kern0pt-ros-args #2}
  \end{tabularx}
\end{tcolorbox}}
% \expected[<note>]{<name>}: the paths (figures/run_<name>.tikz) and the monitor (figures/term_<name>.txt)
\newtcolorbox{seebox}{base, colback=ohmgreen!3, borderline west={2.5pt}{0pt}{ohmgreen!60},
        title={\faEye[regular]~What you should see}, coltitle=ohmgreen!80!black, colbacktitle=ohmgreen!3,
        fonttitle=\small\bfseries, toptitle=2pt, bottomtitle=0pt, breakable=false}
\newcommand{\expected}[2][]{\begin{seebox}
  \centering\resizebox{0.46\linewidth}{!}{\input{figures/run_#2.tikz}}\par\smallskip
  \raggedright{\footnotesize\color{ohmgrey}#1the first lap after the start of the node: white the truth,
    blue the odometry, red ICP. Terminal 1, the monitor (the numbers change from run to run, the ms with the
    computer):}
  \lstinputlisting[basicstyle=\ttfamily\scriptsize, aboveskip=2pt, belowskip=0pt, xleftmargin=4pt,
                   framexleftmargin=3pt, breakindent=1.5em]{figures/term_#2.txt}
\end{seebox}}

\setlist[enumerate,1]{leftmargin=*,itemsep=3pt,topsep=2pt}
\setlist[itemize,1]{leftmargin=*,itemsep=2pt,topsep=2pt}

\begin{document}
\sloppy

% ---------------------------------------------------------------- title
\begin{tcolorbox}[enhanced, colback=ohmred, colframe=ohmred, arc=0pt, left=14pt, right=14pt,
                  top=12pt, bottom=12pt, fontupper=\color{white}]
  {\small Intelligent Robotics (ASY2) · Lab 4}\\[6pt]
  {\LARGE\bfseries ICP}\\[4pt]
  {\large How far did the robot go? Scan matching, alone and with odometry}
\end{tcolorbox}

\noindent
\begin{tabular}{@{}>{\color{ohmgrey}}l l@{\hspace{2.4em}}>{\color{ohmgrey}}l l@{}}
  \faClock[regular]~Duration & 180 minutes & \faUsers~Work in & pairs \\
  \faBook~Script & Cycle 4, section 4.1 & \faFolderOpen[regular]~Files & \datei{ohm\_icp/student/}
\end{tabular}

\section{Overview}
A robot drives a loop around a rack in the hall \thema{production}, lap after lap (\cref{fig:hall}). You do
not steer it: a driver node does. Your node has to find out how the robot moved. It has two sensors:
\begin{itemize}
  \item \textbf{Laser scanner}: 20 scans per second, 360 beams, up to 8\,m. Each range scatters by 1.5\,cm.
  \item \textbf{Odometry}: 50 messages per second from the wheels. Its wheel constants are a little wrong:
        every path is 3\,\% too long, every turn 3\,\% too large.
\end{itemize}
Two scans taken a moment apart see the same racks and walls, a little shifted and turned. The transform that
maps the new scan onto the old one is the motion of the robot in between. ICP finds it. Chain all these
motions and you have the pose of the robot: odometry from the laser.
\begin{itemize}
  \item \textbf{Task 1}: ICP alone, after Besl and McKay. Each match starts at ``the robot did not move''.
  \item \textbf{Task 2}: ICP with the odometry: each match starts where the odometry says the robot went.
\end{itemize}

\begin{figure}[ht]
  \centering
  % generated by tools/make_figures.py — the hall production of the simulator, 1 m = 0.62 cm
\begin{tikzpicture}[x=0.62cm, y=0.62cm, font=\footnotesize]
  \fill[black!4] (0,0) rectangle (20.0,12.0);
  \draw[black!10] (0,0) grid[step=1] (20.0,12.0);
  \fill[black!55] (0.0,11.5) rectangle (20.0,12.0);
  \fill[black!55] (0.0,0.0) rectangle (20.0,0.5);
  \fill[black!55] (0.0,0.5) rectangle (0.5,11.5);
  \fill[black!55] (19.5,0.5) rectangle (20.0,11.5);
  \fill[black!55] (2.5,8.0) rectangle (6.0,10.0);
  \fill[black!55] (8.0,8.0) rectangle (11.5,10.0);
  \fill[black!55] (8.0,3.5) rectangle (11.5,5.5);
  \fill[black!55] (13.5,8.0) rectangle (17.0,10.0);
  \fill[black!55] (13.5,3.5) rectangle (17.0,5.5);
  \fill[black!55] (16.5,1.0) rectangle (18.5,2.5);
  \draw[ohmred, line width=1.1pt] (8.45,2.50) -- (8.65,2.50) -- (8.85,2.50) -- (9.05,2.50) -- (9.25,2.50) -- (9.45,2.50) -- (9.65,2.50) -- (9.85,2.50) -- (10.05,2.50) -- (10.25,2.50) -- (10.45,2.50) -- (10.65,2.50) -- (10.85,2.50) -- (11.05,2.50) -- (11.25,2.50) -- (11.45,2.50) -- (11.65,2.50) -- (11.85,2.51) -- (12.04,2.56) -- (12.21,2.65) -- (12.35,2.79) -- (12.44,2.96) -- (12.49,3.15) -- (12.50,3.35) -- (12.50,3.54) -- (12.50,3.74) -- (12.50,3.94) -- (12.50,4.13) -- (12.50,4.33) -- (12.50,4.53) -- (12.50,4.72) -- (12.50,4.92) -- (12.50,5.12) -- (12.50,5.31) -- (12.50,5.51) -- (12.50,5.71) -- (12.50,5.90) -- (12.49,6.10) -- (12.44,6.29) -- (12.35,6.46) -- (12.21,6.60) -- (12.04,6.69) -- (11.85,6.74) -- (11.65,6.75) -- (11.45,6.75) -- (11.25,6.75) -- (11.05,6.75) -- (10.85,6.75) -- (10.65,6.75) -- (10.45,6.75) -- (10.25,6.75) -- (10.05,6.75) -- (9.85,6.75) -- (9.65,6.75) -- (9.45,6.75) -- (9.25,6.75) -- (9.05,6.75) -- (8.85,6.75) -- (8.65,6.75) -- (8.45,6.75) -- (8.25,6.75) -- (8.05,6.75) -- (7.85,6.75) -- (7.65,6.75) -- (7.45,6.75) -- (7.25,6.75) -- (7.06,6.72) -- (6.88,6.65) -- (6.72,6.53) -- (6.60,6.38) -- (6.53,6.19) -- (6.50,6.00) -- (6.50,5.80) -- (6.50,5.61) -- (6.50,5.41) -- (6.50,5.21) -- (6.50,5.02) -- (6.50,4.82) -- (6.50,4.62) -- (6.50,4.43) -- (6.50,4.23) -- (6.50,4.04) -- (6.50,3.84) -- (6.50,3.64) -- (6.50,3.45) -- (6.50,3.25) -- (6.53,3.06) -- (6.60,2.88) -- (6.72,2.72) -- (6.88,2.60) -- (7.06,2.53) -- (7.25,2.50) -- (7.45,2.50) -- (7.65,2.50) -- (7.85,2.50) -- (8.05,2.50) -- (8.25,2.50) -- cycle;
  \draw[-{Stealth[length=5pt]}, ohmred, line width=1.4pt] (8.5,2.5) -- (9.9,2.5);
  \node[below, ohmred] at (8.5,2.4) {start};
  \draw[-{Stealth[length=4pt]}, line width=0.8pt] (1.2,1.2) -- (3.2,1.2) node[right] {$x$};
  \draw[-{Stealth[length=4pt]}, line width=0.8pt] (1.2,1.2) -- (1.2,3.2) node[above] {$y$};
  \node[below left] at (20.0,0) {20\,m};
  \node[above right] at (0,12.0) {12\,m};
\end{tikzpicture}

  \caption{The hall \thema{production} of the simulator from above, 20\,m \(\times\) 12\,m. Dark: walls and
    racks. The robot drives the red loop counter-clockwise around the lower middle rack at 0.5\,m/s:
    19\,m and 360\,\textdegree{} per lap, about 40\,s.}
  \label{fig:hall}
\end{figure}

\begin{table}[ht]\centering
\caption{Schedule.}
\label{tab:schedule}
\begin{tabular}{@{}rll@{}}
\toprule
min & task & done when \\
\midrule
0--15 & \cref{sec:task0} The scan & \cref{tab:scan} is filled in \\
15--55 & \cref{sec:task1} ICP on two point sets & \thema{pytest -k task1} passes up to 1.5 \\
55--75 & \cref{sec:task1b} ICP in the loop & \cref{tab:alone} is filled in \\
75--90 & \cref{sec:task1c} Pairs that do not belong & \cref{tab:dmax} is filled in \\
90--110 & \cref{sec:task1d} How far apart? & \cref{tab:every} is filled in \\
110--120 & break & \\
120--165 & \cref{sec:task2} ICP with odometry & \cref{tab:odom} is filled in \\
165--180 & \cref{sec:extra} Extra task: the kd-tree & \\
& \cref{sec:extra2} Extra task: every scan & if time is left \\
\bottomrule
\end{tabular}
\end{table}

\paragraph{Notation} as in the script. In the code the names are short:
\begin{center}\small
\begin{tabular}{@{}lll@{}}
\toprule
script & code & \\
\midrule
\(M = \{\vec m_i\}\), \(D = \{\vec d_j\}\) & \thema{M}, \thema{D} & model (the older scan) and data (the newer one), one point per row \\
\(\vec c_m\), \(\vec c_d\) & \thema{c\_m}, \thema{c\_d} & the centroids \\
\(M'\), \(D'\) & \thema{Mc}, \thema{Dc} & both sets shifted by their centroid \\
\(\phi\), \(\vec t\), \(\matsym{T}\) & \thema{phi}, \thema{t}, \thema{T} & rotation, translation, the \(3\times3\) transform \\
\(\omega_{i,j}\) & \thema{pairs\_kept} & the weight of a pair, 1 or 0 \\
\(\epsilon\) & \thema{eps} & the stop: the error changed by less than this \\
\bottomrule
\end{tabular}
\end{center}
Everything is in metres and radians. A pose is a \(3\times3\) transform as \(\matsym{T}\): it maps a point
from the frame of the robot into the hall (\thema{map}).

\newpage
\section{Getting started}
\paragraph{Your own copy.} Fork \datei{ohm\_icp} on GitHub (button \emph{Fork} at the top right), then clone
your fork into the \datei{src/} folder of your workspace. Update the simulator:
\begin{lstlisting}
git clone git@github.com:<your-name>/ohm_icp.git
\end{lstlisting}
\begin{lstlisting}
git -C ir-mecanum-lab pull
\end{lstlisting}
Commit after every task and push, so nothing is lost at the end of the session.

\paragraph{Build.} Build once, then source the workspace in every new terminal:
\begin{lstlisting}
colcon build --symlink-install --packages-select mecanum_lab_interfaces mecanum_lab ohm_icp
\end{lstlisting}
\begin{lstlisting}
source install/setup.bash
\end{lstlisting}
You do not need to build again after a change to your code. Just start the node again.

\paragraph{Simulator.} One launch file starts the simulator in the hall \thema{production}, the driver, the
monitor and RViz. The robot starts to drive at once.
\begin{lstlisting}
ros2 launch ohm_icp icp.launch.py
\end{lstlisting}
The driver steers by the exact pose on \thema{/muster/truth}. Your node never reads \thema{/truth}: it has
to find the motion from the scans (and the odometry) alone. The monitor compares with the truth.

\paragraph{Your files.} \datei{student/icp.py} (the algorithm) and \datei{student/scan\_odometry.py} (from
the messages to the points, from the matches to the pose). You write only the few lines in the middle, marked
\begin{lstlisting}
# the two centroids, and both sets shifted by their own centroid (eq. 4.5 to 4.8)
# TODO STUDENT 1.3
raise NotImplementedError("STUDENT 1.3")
\end{lstlisting}
The comment above it says in words what goes there; the equation is in the task of this handout. Replace the
\thema{raise} line by your code.

\begin{note}
\textbf{numpy} for the point sets: a set of \(N\) points is an \(N\times2\) array, one point per row.
\thema{M[:, 0]} is the column of all \(x\), \thema{M[:, 1]} of all \(y\); \thema{M.mean(axis=0)} the mean
of every column; \thema{M - c} subtracts \thema{c} from every row; \thema{a * b} multiplies element by
element, \thema{np.sum(a)} adds up; \thema{A @ B} is the matrix product; \thema{np.linalg.inv(A)} the
inverse.
\end{note}

\paragraph{The node.} \datei{student/icp\_node.py} is given. It subscribes to the scan and the odometry,
calls your functions and publishes the pose. It runs in a second terminal:
\begin{lstlisting}
ros2 run ohm_icp icp_node.py --ros-args -p task:=1
\end{lstlisting}
If a function is not written yet, the node stops with \thema{NotImplementedError: STUDENT 1.1}: that is the
next one to write.

\paragraph{Self check.} The tests try your functions on made-up point sets, one task at a time:
\begin{lstlisting}
python3 -m pytest test/test_icp.py -q -k task1
\end{lstlisting}

\paragraph{Three terminals.} All of them in the workspace, sourced.
\begin{center}\small
\begin{tabular}{@{}>{\color{ohmgrey}}l l l@{}}
  terminal 1 & simulator, driver, monitor, RViz & the monitor prints here every 10\,s \\
  terminal 2 & your node           & Ctrl-C and start again after every change \\
  terminal 3 & tests, hints, plots, \thema{ros2 topic} & \\
\end{tabular}
\end{center}
Every task starts with a box \emph{Running}: what terminals 1 and 2 run. It ends with \emph{What you should
see}: the first lap and the monitor when the task is done. RViz uses the same colours: white the true path,
blue the odometry, red your ICP. All three start where your node started.

\paragraph{The monitor and the plots.} Each start of your node is a new run. Every 10\,s the monitor prints,
from the start of the run: how far the robot went (\emph{driven}), how much it turned (\emph{turned}, one lap
is 360\,\textdegree), and the error of the odometry and of ICP against the truth, in position and heading.
After each lap it prints the errors again. The monitor writes every run to \datei{runs/}; the last one as
plots:
\begin{lstlisting}
ros2 run ohm_icp plot_run.py
\end{lstlisting}

\paragraph{Hints.} Stuck? Every block has three hints: the idea, how to write it in Python, and at last
the code. Each call opens the next one, in the folder \datei{ohm\_icp}:
\begin{lstlisting}
python3 hint.py 1.3
\end{lstlisting}
\thema{python3 hint.py} alone shows the hints you have opened. Try on your own first: the hints notice
when you ask too early.

\begin{figure}[!hb]
  \centering
  \resizebox{\linewidth}{!}{% The ROS graph of the lab. Red: your node. Grey: given.
\begin{tikzpicture}[font=\footnotesize, >={Stealth[length=5pt]},
    node/.style={rectangle, rounded corners=2pt, draw=black!50, fill=black!4, minimum height=0.9cm,
                 minimum width=2.4cm, align=center},
    mine/.style={node, draw=ohmred, fill=ohmred!8, line width=0.9pt},
    topic/.style={font=\scriptsize\ttfamily, fill=white, inner sep=1.5pt}]
  \node[node] (sim) at (0,0) {simulator\\\scriptsize\texttt{mecanum\_sim}};
  \node[mine] (icp) at (6.2,0) {your ICP\\\scriptsize\texttt{icp\_node.py}};
  \node[node] (mon) at (11.6,0) {monitor\\\scriptsize\texttt{monitor\_node.py}};
  \node[node] (drv) at (0,-2.4) {driver\\\scriptsize\texttt{route\_node.py}};
  \node[node] (rviz) at (11.6,-2.4) {RViz};
  \draw[->] (sim.20) -- node[topic, above=1pt] {/muster/scan} (icp.160);
  \draw[->] (sim.-20) -- node[topic, below=1pt] {/muster/odom} (icp.200);
  \draw[->] (icp) -- node[topic, above=1pt] {/muster/icp/info} (mon);
  \draw[->] (sim.250) -- node[topic, left=1pt] {/muster/truth} (drv.110);
  \draw[->] (drv.70) -- node[topic, right=1pt] {/muster/cmd\_vel} (sim.290);
  \draw[->, black!50] (sim.90) -- ++(0,0.7) -| node[topic, pos=0.25, above=1pt] {/muster/truth, /muster/odom}
        (mon.90);
  \draw[->] (mon) -- node[topic, right=1pt] {/muster/icp\_lab/*} (rviz);
  \draw[->] (icp.270) |- node[topic, pos=0.75, below=1pt] {/muster/icp/pose, /muster/icp/markers} (rviz.180);
\end{tikzpicture}
}
  \caption{The nodes and topics of the lab. Red: your node. Grey: given.}
  \label{fig:nodes}
\end{figure}

\clearpage
\section{Task 0: the scan\minutes{15}}
\label{sec:task0}
\running{}{(nothing yet)}

Before you match scans, look at one. In terminal 3:
\begin{lstlisting}
ros2 topic hz /muster/scan
\end{lstlisting}
\begin{lstlisting}
ros2 topic echo --once /muster/scan
\end{lstlisting}
\begin{lstlisting}
ros2 interface show sensor_msgs/msg/LaserScan
\end{lstlisting}
\thema{echo} prints a long list of \thema{ranges}. Add \thema{-{}-field angle\_increment} (or another field)
to see one field only. Fill in \cref{tab:scan}.

\begin{table}[ht]\centering
\renewcommand{\arraystretch}{1.5}
\caption{The laser scanner. Fill in.}
\label{tab:scan}
\begin{tabular}{@{}>{\raggedright\arraybackslash}p{7.4cm}p{6.0cm}@{}}
\toprule
& \thema{/muster/scan} \\
\midrule
message type & \\
scans per second & \\
frame of the points (\thema{header.frame\_id}) & \\
\thema{angle\_min}, \thema{angle\_increment} (in degrees) & \\
number of beams & \\
\thema{range\_min}, \thema{range\_max} & \\
the range of a beam that hits nothing & \\
\bottomrule
\end{tabular}
\end{table}

\begin{itemize}
  \item Beam 0 points forward (\(x\) of the robot). In which direction does beam 90 point?
  \item Watch the simulator window: the scan is drawn as dots. Which beams hit nothing on the loop?
\end{itemize}

\begin{checkpoint}
A beam that hits nothing is not missing: it reads \thema{range\_max}. It is no point of a wall, and your code
has to drop it.
\end{checkpoint}

\newpage
\section{Task 1: ICP on two point sets\minutes{40}}
\label{sec:task1}
\running{}{(nothing yet, the tests first)}

First the algorithm itself, without ROS: \datei{icp.py}, the four steps of the script.

\begin{script}{section 4.1.3, ICP after Besl and McKay}
\begin{enumerate}[label=\arabic*.]
  \item For each point \(\vec d_j \in D\) find the closest point \(\vec m_i \in M\).
  \item From these pairs compute \(\matsym R\) and \(\vec t\) that minimise
        \begin{equation}
          E(\matsym R, \vec t) = \frac{1}{N} \sum_{i=1}^{N} \|\vec m_i - (\matsym R \vec d_i + \vec t)\|^2 .
          \tag{4.4}
        \end{equation}
  \item Apply the transform to \(D\).
  \item If the error changed by less than \(\epsilon\), stop; otherwise go to 1.
\end{enumerate}
Step 2 in closed form. The centroids, and both sets shifted by them:
\begin{equation}
  \vec c_m = \frac1N \sum_{i=1}^{N} \vec m_i, \quad \vec c_d = \frac1N \sum_{i=1}^{N} \vec d_i, \qquad
  \vec m'_i = \vec m_i - \vec c_m, \quad \vec d'_i = \vec d_i - \vec c_d .
  \tag{4.5--4.8}
\end{equation}
The rotation and the translation:
\begin{align}
  \phi &= \operatorname{atan2}\Big( \sum_{i=1}^{N} \big(y'_{m,i}\, x'_{d,i} - x'_{m,i}\, y'_{d,i}\big),\;
                                    \sum_{i=1}^{N} \big(x'_{m,i}\, x'_{d,i} + y'_{m,i}\, y'_{d,i}\big) \Big) ,
  \tag{4.14}\\
  \vec t &= \vec c_m - \matsym R(\phi)\, \vec c_d ,
  \tag{4.17}
\end{align}
and as one transform
\(\matsym T = \big(\begin{smallmatrix} \cos\phi & -\sin\phi & t_x\\ \sin\phi & \cos\phi & t_y\\ 0&0&1
\end{smallmatrix}\big)\) (4.18).
\end{script}

\begin{todo}{1.2}
\thema{closest\_points(D, M)}, step 1, brute force: for every row of \thema{D} the index of the closest row of
\thema{M} and its distance. \thema{diff}, all differences \(\vec d_j - \vec m_i\), is given.
\end{todo}

\begin{todo}{1.3}
\thema{best\_transform(M, D)}: the centroids and the shifted sets (4.5--4.8). Row \(i\) of \thema{M} and of
\thema{D} is pair \(i\).
\end{todo}

\begin{todo}{1.4}
\thema{best\_transform}: \(\phi\) with (4.14) and \(\vec t\) with (4.17). \thema{rotation(phi)} and
\thema{se2(x, y, phi)} are given.
\end{todo}

\begin{todo}{1.5}
\thema{icp(M, D)}, step 3: move \thema{D\_now} by the step \thema{T\_step}, and put \thema{T\_step} on top of
the transform so far, \thema{T}. Steps 1, 2 and 4 are given around it.
\end{todo}

\begin{note}
The order of the product. ICP finds \(\matsym T_1\) in the first iteration, applies it, then finds
\(\matsym T_2\) for the points that already moved. A point of \(D\) ends at
\(\matsym T_2 \matsym T_1 \vec d\): the newest step stands on the \emph{left}. After \(n\) iterations
\(\matsym T = \matsym T_n \cdots \matsym T_2 \matsym T_1\).
\end{note}

\thema{pytest -k task1} also tries the exercise of the script (\(M = \{(5,4), (6,2)\}\),
\(D = \{(1.5,2.7), (2,0.5)\}\) gives \(\phi = 13.76\)\,\textdegree{} and \(\vec t = (4.18, 1.03)\)) and a
room with a box in it, moved and turned. Tests 1.1, 1.6 and 1.7 fail until the next tasks.

\begin{checkpoint}
\thema{pytest -k "task1 and not task1\_1 and not task1\_6 and not task1\_7"} passes. Why is one iteration
enough for the exercise of the script, and not for the room?
\end{checkpoint}

\newpage
\section{Task 1: ICP in the loop\minutes{20}}
\label{sec:task1b}
\running{}{-p task:=1 -p d\_max:=0}

Now the scans of the robot. The node keeps the last matched scan as the model \(M\). Every tenth scan (twice a
second) becomes the data \(D\): ICP maps it onto the model, the result \(\matsym T\) is the motion of the robot
in between. Then the new scan is the model for the next match.

\begin{todo}{1.1}
\datei{scan\_odometry.py}, \thema{scan\_to\_points(msg)}: the points of one scan in the frame of the laser,
\((r \cos\alpha,\ r \sin\alpha)\) with \(\alpha\) = \thema{angle\_min} + \(i\,\cdot\) \thema{angle\_increment}.
Drop every beam without an echo (\cref{tab:scan}).
\end{todo}

\(\matsym T\) maps the new scan onto the old one: it is the pose of the robot now, in the frame of the robot at
the old scan. The pose at the old scan maps that frame into the hall.

\begin{todo}{1.6}
\thema{next\_pose(pose, T)}: the pose at the new scan, from the pose at the old scan and \(\matsym T\). Both
are \(3\times3\) transforms.
\end{todo}

Start the node. It starts where the odometry is. Watch RViz: grey the model, amber the new scan where ICP
starts, red the new scan where ICP put it, light blue the pairs. Let the robot drive one full lap
(\(\approx 40\)\,s) and take the line \thema{lap 1} of the monitor. Fill in \cref{tab:alone}.

\begin{table}[ht]\centering
\renewcommand{\arraystretch}{1.4}
\caption{One lap, ICP alone with \thema{d\_max:=0}. Fill in.}
\label{tab:alone}
\begin{tabular}{@{}p{3.0cm}p{2.3cm}p{2.3cm}p{2.3cm}p{2.3cm}@{}}
\toprule
& driven in m & turned in \textdegree & error in m & error in \textdegree \\
\midrule
truth & & 360 & --- & --- \\
odometry & & & & \\
ICP & & & & \\
\bottomrule
\end{tabular}
\end{table}

\begin{itemize}
  \item Which of the two is better in heading? Which one in distance?
  \item In RViz, watch the pairs while the robot passes a rack. Do all of them connect the same wall?
\end{itemize}

\expected[\thema{d\_max:=0}; ]{dmax0}

\begin{checkpoint}
The ICP path has about the right shape, but it is too short, and its heading is no better than the one of
the odometry. The next task finds out why.
\end{checkpoint}

\newpage
\section{Task 1: pairs that do not belong\minutes{15}}
\label{sec:task1c}
\running{}{-p task:=1 \textsf{, then} -p d\_max:=1.0}

The new scan sees things the old one did not: the side of a rack that came into view, a wall that was out of
range. Step 1 still finds a closest point for each of them, a wrong pair. Wrong pairs pull the result off. In
(4.3) of the script every pair has a weight \(\omega_{i,j}\), 1 or 0. Give the weight 0 to every pair farther
apart than \thema{d\_max}.

\begin{todo}{1.7}
\datei{icp.py}, \thema{pairs\_kept(distance, d\_max)}: \thema{True} for every pair whose distance is at most
\thema{d\_max}.
\end{todo}

The default is \thema{d\_max} = 0.3\,m. Run one lap with it, then one with 1.0\,m. Fill in \cref{tab:dmax}
(the row 0 is from \cref{tab:alone}).

\begin{table}[ht]\centering
\renewcommand{\arraystretch}{1.4}
\caption{One lap, ICP alone, three values of \thema{d\_max}. Fill in.}
\label{tab:dmax}
\begin{tabular}{@{}p{3.0cm}p{2.3cm}p{2.3cm}p{2.3cm}p{2.3cm}@{}}
\toprule
\thema{d\_max} in m & driven in m & turned in \textdegree & error in m & error in \textdegree \\
\midrule
0 (all pairs) & & & & \\
0.3 & & & & \\
1.0 & & & & \\
\bottomrule
\end{tabular}
\end{table}

\begin{itemize}
  \item Why does 1.0\,m help less than 0.3\,m?
  \item The robot moves 25\,cm between two matches on a straight. What would happen with \thema{d\_max} = 0.1\,m?
        Try it.
\end{itemize}

\expected[\thema{d\_max:=0.3}; ]{alone}

\begin{checkpoint}
\thema{pytest -k task1} passes. With \thema{d\_max} = 0.3\,m ICP alone is better than the odometry in
heading by more than a factor two. In position it is not better: its path is still too short.
\end{checkpoint}

\newpage
\section{Task 1: how far apart?\minutes{20}}
\label{sec:task1d}
\running{}{-p task:=1 -p every:=2 \textsf{, then 5 and 20}}

\thema{every} says which scans are matched: 10 is every tenth scan, twice a second. Try 2, 5 and 20, one lap
each (\cref{tab:every}, the row 10 is from \cref{tab:dmax}).

\begin{table}[ht]\centering
\renewcommand{\arraystretch}{1.4}
\caption{One lap, ICP alone, \thema{d\_max} = 0.3\,m. Fill in. ms: per match, terminal 2.}
\label{tab:every}
\begin{tabular}{@{}p{1.9cm}p{2.1cm}p{2.1cm}p{2.1cm}p{2.1cm}p{2.1cm}@{}}
\toprule
\thema{every} & driven in m & turned in \textdegree & error in m & error in \textdegree & ms \\
\midrule
2 & & & & & \\
5 & & & & & \\
10 & & & & & \\
20 & & & & & \\
\bottomrule
\end{tabular}
\end{table}

\begin{itemize}
  \item \thema{every} = 2: the path is far too short. Each match finds a little too little motion: the beams
        of both scans hit the wall at almost the same places, and ICP stops where every point sits on its old
        neighbour, a local minimum. What does that small error per match do when there are five times as many
        matches?
  \item \thema{every} = 20: half a metre and up to 36\,\textdegree{} in a corner between two scans. Look at the
        first match after a corner in RViz (\cref{fig:match}, left). Why does ICP fall into the wrong minimum
        here? What does \thema{d\_max} do to the right pairs, which are now more than 0.3\,m apart?
\end{itemize}

\begin{figure}[ht]
  \centering
  \resizebox{0.92\linewidth}{!}{% generated by tools/make_figures.py: one match 20 scans apart in a corner, in the frame of the model
\begin{tikzpicture}[x=0.38cm, y=0.38cm]
  \begin{scope}\clip (0.00,0.00) rectangle (9.00,9.00);
  \fill[black!85] (0.00,0.00) rectangle (9.00,9.00);
  \foreach \x/\y in {1.72/0.00, 1.79/0.03, 1.81/0.06, 1.80/0.09, 1.85/0.13, 1.87/0.16, 1.90/0.20, 1.92/0.24, 1.99/0.28, 2.00/0.32, 2.03/0.36, 2.06/0.40, 2.13/0.45, 2.15/0.50, 2.20/0.55, 2.25/0.60, 2.31/0.66, 2.32/0.71, 2.38/0.77, 2.41/0.83, 2.48/0.90, 2.55/0.98, 2.61/1.06, 2.69/1.14, 2.77/1.23, 4.43/3.10, 4.40/3.20, 3.05/4.69, 2.80/4.48, 2.73/4.55, 2.67/4.62, 2.59/4.68, 2.52/4.75, 2.44/4.79, 1.23/2.52, 1.11/2.37, 1.00/2.24, 0.90/2.13, 0.81/2.00, 0.74/1.92, 0.67/1.84, 0.61/1.79, 0.55/1.70, 0.50/1.64, 0.46/1.59, 0.42/1.55, 0.38/1.51, 0.33/1.43, 0.29/1.39, 0.26/1.34, 0.23/1.31, 0.20/1.28, 0.18/1.25, 0.15/1.24, 0.13/1.20, 0.10/1.14, 0.08/1.16, 0.06/1.11, 0.04/1.09, 0.02/1.06, 0.00/1.08, -0.02/1.00, -0.03/1.00, -0.05/1.00, -0.07/1.04, -0.09/1.02, -0.11/1.05, -0.13/1.04, -0.16/1.11, -0.18/1.12, -0.20/1.13, -0.22/1.13, -0.25/1.18, -0.27/1.16, -0.31/1.24, -0.32/1.21, -0.36/1.25, -0.39/1.27, -0.43/1.31, -0.47/1.35, -0.50/1.37, -0.55/1.42, -0.58/1.43, -0.63/1.48, -0.67/1.50, -0.72/1.55, -0.80/1.63, -0.86/1.68, -0.92/1.74, -0.99/1.78, -1.07/1.85, -1.14/1.90, -1.26/2.02, -1.38/2.12, -1.47/2.18, -1.61/2.31, -1.76/2.43, -1.95/2.58, -2.17/2.78, -2.39/2.95, -2.69/3.20, -4.50/0.79, -4.36/0.69, -4.22/0.59, -4.16/0.51, -4.04/0.42, -3.96/0.35, -3.84/0.27, -3.78/0.20, -3.70/0.13, -3.62/0.06, -3.51/0.00, -3.45/-0.06, -3.41/-0.12, -3.31/-0.17, -3.25/-0.23, -3.21/-0.28, -3.12/-0.33, -3.09/-0.38, -3.02/-0.42, -2.96/-0.47, -2.92/-0.51, -2.86/-0.56, -2.81/-0.60, -2.76/-0.64, -2.72/-0.68, -2.68/-0.72, -2.66/-0.76, -2.59/-0.79, -2.55/-0.83, -2.49/-0.86, -2.46/-0.90, -2.42/-0.93, -2.36/-0.95, -2.33/-0.99, -2.31/-1.03, -2.28/-1.06, -2.24/-1.09, -2.19/-1.12, -2.16/-1.15, -2.11/-1.17, -2.11/-1.22, -2.04/-1.22, -2.03/-1.27, -1.97/-1.28, -1.95/-1.32, -1.91/-1.34, -1.90/-1.38, -1.86/-1.40, -1.82/-1.42, -1.80/-1.46, -1.76/-1.48, -1.74/-1.51, -1.70/-1.53, -1.67/-1.55, -1.64/-1.58, -1.60/-1.60, -1.58/-1.63, -1.54/-1.65, -1.52/-1.68, -1.49/-1.71, -1.48/-1.76, -1.43/-1.77, -1.40/-1.79, -1.37/-1.81, -1.35/-1.86, -1.31/-1.88, -1.28/-1.89, -1.26/-1.93, -1.22/-1.95, -1.18/-1.96, -1.15/-2.00, -1.11/-2.01, -1.07/-2.01, -1.06/-2.07, -1.02/-2.10, -0.98/-2.11, -0.97/-2.18, -0.92/-2.17, -0.89/-2.20, -0.86/-2.23, -0.82/-2.27, -0.79/-2.30, -0.75/-2.30, -0.72/-2.37, -0.68/-2.36, -0.65/-2.42, -0.61/-2.44, -0.57/-2.48, -0.53/-2.50, -0.49/-2.52, -0.45/-2.57, -0.41/-2.62, -0.37/-2.64, -0.33/-2.70, -0.29/-2.72, -0.24/-2.76, -0.20/-2.79, -0.15/-2.81, -0.10/-2.85, -0.05/-2.91, -0.00/-2.93, 0.05/-2.99, 0.11/-3.01, 0.16/-3.09, 0.22/-3.13, 0.28/-3.18, 0.34/-3.25, 0.40/-3.28, 0.47/-3.37, 0.54/-3.38, 0.61/-3.45, 0.68/-3.50, 0.76/-3.57, 0.84/-3.64, 0.93/-3.72, 1.02/-3.81, 1.11/-3.87, 1.21/-3.96, 2.31/-4.00, 2.34/-3.90, 2.42/-3.87, 2.46/-3.79, 2.52/-3.74, 2.57/-3.68, 2.63/-3.63, 2.69/-3.57, 2.72/-3.48, 2.77/-3.42, 2.84/-3.38, 2.87/-3.30, 2.92/-3.25, 2.98/-3.20, 3.01/-3.12, 3.08/-3.08, 3.13/-3.03, 3.54/-3.30, 4.09/-2.46, 3.71/-2.14, 3.42/-1.89, 3.15/-1.67, 2.92/-1.49, 2.73/-1.33, 2.60/-1.21, 2.46/-1.09, 2.31/-0.98, 2.21/-0.89, 2.12/-0.81, 2.03/-0.74, 1.94/-0.67, 1.87/-0.61, 1.79/-0.55, 1.77/-0.51, 1.69/-0.45, 1.64/-0.41, 1.59/-0.37, 1.54/-0.33, 1.49/-0.29, 1.52/-0.27, 1.52/-0.24, 1.57/-0.22, 1.56/-0.19, 1.62/-0.17, 1.63/-0.14, 1.63/-0.11, 1.68/-0.09, 1.70/-0.06, 1.72/-0.03} \fill[black!25] ({\x - (-4.5) + 0.0},{\y - (-4.0)}) circle (0.9pt);
  \foreach \x/\y in {4.49/2.09, 4.47/2.18, 4.45/2.27, 4.44/2.36, 4.38/2.43, 4.38/2.53, 4.36/2.62, 4.36/2.72, 4.31/2.80, 4.27/2.88, 1.97/1.38, 1.86/1.35, 1.74/1.31, 1.69/1.32, 1.59/1.29, 1.50/1.26, 1.43/1.24, 1.38/1.25, 1.31/1.22, 1.25/1.21, 1.18/1.18, 1.14/1.18, 1.10/1.17, 1.02/1.13, 0.97/1.12, 0.94/1.13, 0.91/1.12, 0.87/1.11, 0.83/1.11, 0.79/1.08, 0.75/1.07, 0.69/1.02, 0.70/1.07, 0.66/1.05, 0.63/1.05, 0.60/1.04, 0.57/1.04, 0.55/1.03, 0.52/1.01, 0.49/1.00, 0.47/1.00, 0.44/0.99, 0.42/1.00, 0.40/0.98, 0.38/0.98, 0.36/0.98, 0.34/0.98, 0.32/0.98, 0.29/0.96, 0.28/0.97, 0.25/0.95, 0.23/0.94, 0.22/0.93, 0.20/0.94, 0.18/0.93, 0.16/0.91, 0.15/0.94, 0.13/0.95, 0.12/1.00, 0.11/1.05, 0.10/1.10, 0.08/1.18, 0.06/1.23, 0.05/1.32, 0.02/1.40, 0.00/1.50, -0.03/1.61, -0.06/1.73, -0.10/1.90, -0.15/2.10, -0.20/2.31, -0.27/2.59, -0.37/2.99, -0.48/3.44, -0.66/4.16, -3.92/4.85, -3.85/4.58, -3.81/4.39, -3.76/4.18, -3.72/3.99, -3.67/3.80, -3.61/3.61, -3.59/3.46, -3.53/3.29, -3.49/3.14, -3.47/3.01, -3.42/2.87, -3.41/2.76, -3.37/2.63, -3.34/2.52, -3.33/2.42, -3.30/2.31, -3.27/2.21, -3.24/2.11, -3.23/2.02, -3.20/1.93, -3.18/1.84, -3.17/1.76, -3.14/1.67, -3.12/1.59, -3.09/1.51, -3.09/1.44, -3.08/1.37, -3.06/1.30, -3.04/1.23, -3.02/1.16, -3.00/1.09, -2.95/1.01, -2.97/0.97, -2.95/0.90, -2.92/0.84, -2.92/0.78, -2.87/0.72, -2.89/0.67, -2.86/0.61, -2.88/0.56, -2.84/0.50, -2.83/0.45, -2.83/0.40, -2.78/0.34, -2.81/0.29, -2.78/0.24, -2.77/0.19, -2.77/0.15, -2.75/0.10, -2.71/0.05, -2.70/0.00, -2.72/-0.05, -2.69/-0.09, -2.70/-0.14, -2.65/-0.19, -2.66/-0.23, -2.66/-0.28, -2.67/-0.33, -2.62/-0.37, -2.61/-0.41, -2.59/-0.46, -2.61/-0.51, -2.59/-0.55, -2.57/-0.59, -2.57/-0.64, -2.55/-0.68, -2.54/-0.73, -2.53/-0.77, -2.51/-0.82, -2.51/-0.86, -2.49/-0.91, -2.47/-0.95, -2.48/-1.00, -2.47/-1.05, -2.45/-1.09, -2.43/-1.13, -2.43/-1.18, -2.42/-1.23, -2.41/-1.28, -2.35/-1.30, -2.40/-1.39, -2.37/-1.43, -2.37/-1.48, -2.35/-1.52, -2.34/-1.58, -2.31/-1.62, -2.31/-1.68, -2.31/-1.74, -2.28/-1.78, -2.26/-1.83, -2.23/-1.88, -2.25/-1.95, -2.23/-2.01, -2.20/-2.06, -2.19/-2.12, -2.19/-2.19, -2.16/-2.24, -2.16/-2.31, -2.13/-2.37, -2.11/-2.42, -2.10/-2.51, -2.07/-2.55, -2.07/-2.65, -2.03/-2.70, -2.03/-2.79, -2.02/-2.88, -1.97/-2.93, -1.97/-3.04, -1.95/-3.13, -1.93/-3.21, -1.91/-3.31, -1.88/-3.40, -1.86/-3.49, -1.83/-3.59, -1.80/-3.68, -1.78/-3.81, -1.74/-3.90, 1.31/-3.81, 1.16/-3.19, 1.05/-2.73, 0.95/-2.35, 0.89/-2.10, 0.82/-1.85, 0.80/-1.71, 0.74/-1.53, 0.73/-1.43, 0.70/-1.31, 0.67/-1.20, 0.64/-1.10, 0.63/-1.05, 0.63/-1.01, 0.64/-0.98, 0.67/-1.00, 0.70/-1.00, 0.73/-1.00, 0.74/-0.98, 0.76/-0.97, 0.78/-0.96, 0.81/-0.97, 0.83/-0.95, 0.88/-0.97, 0.90/-0.96, 0.92/-0.96, 0.94/-0.94, 0.98/-0.94, 0.97/-0.91, 1.03/-0.93, 1.06/-0.92, 1.07/-0.90, 1.10/-0.89, 1.13/-0.88, 1.17/-0.88, 1.21/-0.88, 1.22/-0.86, 1.27/-0.86, 1.29/-0.84, 1.33/-0.83, 1.38/-0.83, 1.42/-0.82, 1.43/-0.79, 1.52/-0.81, 1.56/-0.80, 1.60/-0.78, 1.66/-0.77, 1.68/-0.75, 1.75/-0.74, 1.81/-0.73, 1.85/-0.71, 1.93/-0.70, 2.00/-0.69, 2.06/-0.67, 2.09/-0.64, 2.19/-0.63, 2.27/-0.61, 2.35/-0.58, 2.44/-0.56, 2.55/-0.54} \fill[ohmamber] ({\x - (-4.5) + 0.0},{\y - (-4.0)}) circle (0.8pt);
  \foreach \x/\y in {4.40/4.27, 3.85/3.90, 3.41/3.61, 3.35/3.67, 3.30/3.75, 3.25/3.82, 3.21/3.91, 3.13/3.95, 3.09/4.04, 3.05/4.12, 3.01/4.22, 2.94/4.27, 2.87/4.34, 1.23/2.14, 1.14/2.07, 1.04/1.99, 0.98/1.98, 0.90/1.92, 0.83/1.86, 0.77/1.82, 0.72/1.81, 0.66/1.76, 0.62/1.73, 0.56/1.68, 0.52/1.67, 0.48/1.64, 0.42/1.58, 0.38/1.55, 0.35/1.54, 0.32/1.53, 0.29/1.50, 0.26/1.49, 0.22/1.45, 0.19/1.43, 0.15/1.36, 0.14/1.41, 0.11/1.38, 0.09/1.37, 0.06/1.35, 0.04/1.33, 0.01/1.32, -0.01/1.29, -0.03/1.27, -0.05/1.26, -0.07/1.24, -0.09/1.24, -0.11/1.22, -0.13/1.21, -0.15/1.21, -0.17/1.20, -0.18/1.19, -0.20/1.16, -0.22/1.17, -0.23/1.14, -0.25/1.13, -0.26/1.11, -0.28/1.11, -0.30/1.10, -0.31/1.08, -0.33/1.10, -0.35/1.10, -0.37/1.15, -0.40/1.19, -0.43/1.22, -0.47/1.30, -0.51/1.34, -0.56/1.42, -0.60/1.49, -0.66/1.57, -0.73/1.67, -0.80/1.77, -0.89/1.91, -1.01/2.09, -1.13/2.26, -1.29/2.50, -1.52/2.85, -1.78/3.23, -2.19/3.84, -4.44/1.80, -4.35/1.68, -4.30/1.58, -4.21/1.47, -4.15/1.38, -4.10/1.29, -4.04/1.20, -3.98/1.11, -3.91/1.02, -3.87/0.95, -3.82/0.87, -3.77/0.79, -3.73/0.72, -3.67/0.65, -3.62/0.58, -3.56/0.52, -3.54/0.45, -3.51/0.39, -3.47/0.33, -3.42/0.27, -3.39/0.21, -3.34/0.16, -3.26/0.10, -3.27/0.05, -3.23/-0.00, -3.18/-0.06, -3.15/-0.11, -3.09/-0.15, -3.09/-0.20, -3.05/-0.25, -3.04/-0.30, -2.98/-0.34, -2.96/-0.39, -2.94/-0.44, -2.87/-0.47, -2.88/-0.53, -2.84/-0.56, -2.82/-0.61, -2.79/-0.65, -2.76/-0.69, -2.70/-0.72, -2.68/-0.77, -2.68/-0.82, -2.64/-0.85, -2.63/-0.90, -2.57/-0.92, -2.56/-0.97, -2.55/-1.02, -2.54/-1.06, -2.48/-1.08, -2.45/-1.12, -2.42/-1.16, -2.42/-1.21, -2.39/-1.24, -2.36/-1.28, -2.34/-1.32, -2.31/-1.36, -2.28/-1.39, -2.25/-1.43, -2.22/-1.47, -2.20/-1.51, -2.17/-1.55, -2.14/-1.58, -2.13/-1.63, -2.10/-1.67, -2.07/-1.71, -2.04/-1.74, -2.02/-1.78, -1.99/-1.82, -1.97/-1.87, -1.91/-1.87, -1.92/-1.96, -1.88/-1.99, -1.86/-2.04, -1.83/-2.08, -1.80/-2.13, -1.76/-2.15, -1.74/-2.21, -1.72/-2.27, -1.67/-2.29, -1.64/-2.33, -1.60/-2.37, -1.58/-2.44, -1.55/-2.49, -1.51/-2.53, -1.48/-2.58, -1.45/-2.64, -1.41/-2.68, -1.37/-2.75, -1.33/-2.79, -1.29/-2.84, -1.26/-2.92, -1.21/-2.95, -1.18/-3.03, -1.13/-3.07, -1.09/-3.16, -1.05/-3.23, -0.99/-3.26, -0.95/-3.37, -0.90/-3.44, -0.85/-3.51, -0.80/-3.60, -0.75/-3.68, -0.69/-3.76, -0.63/-3.83, -0.57/-3.91, 1.12/-3.96, 1.19/-3.92, 1.25/-3.86, 1.32/-3.83, 1.38/-3.80, 1.45/-3.76, 1.51/-3.70, 1.59/-3.71, 1.64/-3.64, 2.40/-2.97, 2.04/-2.43, 1.78/-2.04, 1.55/-1.71, 1.41/-1.50, 1.27/-1.29, 1.19/-1.17, 1.08/-1.01, 1.03/-0.93, 0.96/-0.83, 0.90/-0.74, 0.83/-0.65, 0.81/-0.61, 0.79/-0.57, 0.79/-0.54, 0.83/-0.54, 0.86/-0.54, 0.88/-0.53, 0.89/-0.50, 0.90/-0.49, 0.92/-0.47, 0.95/-0.47, 0.96/-0.45, 1.01/-0.45, 1.03/-0.43, 1.05/-0.42, 1.06/-0.40, 1.09/-0.39, 1.08/-0.35, 1.14/-0.35, 1.17/-0.34, 1.17/-0.31, 1.19/-0.29, 1.22/-0.28, 1.25/-0.26, 1.29/-0.25, 1.30/-0.22, 1.34/-0.20, 1.35/-0.18, 1.39/-0.16, 1.43/-0.14, 1.47/-0.12, 1.47/-0.09, 1.56/-0.07, 1.59/-0.05, 1.62/-0.02, 1.68/0.01, 1.69/0.04, 1.76/0.07, 1.80/0.10, 1.84/0.13, 1.91/0.17, 1.97/0.20, 2.02/0.24, 2.03/0.28, 2.13/0.33, 2.19/0.37, 2.26/0.42, 2.34/0.47, 2.44/0.53} \fill[ohmred] ({\x - (-4.5) + 0.0},{\y - (-4.0)}) circle (0.8pt);
  \fill[white] (4.50,4.00) circle (2pt);
  \node[below right, white, font=\scriptsize] at (0.00,9.00) {task 1};
  \end{scope}
  \begin{scope}\clip (10.00,0.00) rectangle (19.00,9.00);
  \fill[black!85] (10.00,0.00) rectangle (19.00,9.00);
  \foreach \x/\y in {1.72/0.00, 1.79/0.03, 1.81/0.06, 1.80/0.09, 1.85/0.13, 1.87/0.16, 1.90/0.20, 1.92/0.24, 1.99/0.28, 2.00/0.32, 2.03/0.36, 2.06/0.40, 2.13/0.45, 2.15/0.50, 2.20/0.55, 2.25/0.60, 2.31/0.66, 2.32/0.71, 2.38/0.77, 2.41/0.83, 2.48/0.90, 2.55/0.98, 2.61/1.06, 2.69/1.14, 2.77/1.23, 4.43/3.10, 4.40/3.20, 3.05/4.69, 2.80/4.48, 2.73/4.55, 2.67/4.62, 2.59/4.68, 2.52/4.75, 2.44/4.79, 1.23/2.52, 1.11/2.37, 1.00/2.24, 0.90/2.13, 0.81/2.00, 0.74/1.92, 0.67/1.84, 0.61/1.79, 0.55/1.70, 0.50/1.64, 0.46/1.59, 0.42/1.55, 0.38/1.51, 0.33/1.43, 0.29/1.39, 0.26/1.34, 0.23/1.31, 0.20/1.28, 0.18/1.25, 0.15/1.24, 0.13/1.20, 0.10/1.14, 0.08/1.16, 0.06/1.11, 0.04/1.09, 0.02/1.06, 0.00/1.08, -0.02/1.00, -0.03/1.00, -0.05/1.00, -0.07/1.04, -0.09/1.02, -0.11/1.05, -0.13/1.04, -0.16/1.11, -0.18/1.12, -0.20/1.13, -0.22/1.13, -0.25/1.18, -0.27/1.16, -0.31/1.24, -0.32/1.21, -0.36/1.25, -0.39/1.27, -0.43/1.31, -0.47/1.35, -0.50/1.37, -0.55/1.42, -0.58/1.43, -0.63/1.48, -0.67/1.50, -0.72/1.55, -0.80/1.63, -0.86/1.68, -0.92/1.74, -0.99/1.78, -1.07/1.85, -1.14/1.90, -1.26/2.02, -1.38/2.12, -1.47/2.18, -1.61/2.31, -1.76/2.43, -1.95/2.58, -2.17/2.78, -2.39/2.95, -2.69/3.20, -4.50/0.79, -4.36/0.69, -4.22/0.59, -4.16/0.51, -4.04/0.42, -3.96/0.35, -3.84/0.27, -3.78/0.20, -3.70/0.13, -3.62/0.06, -3.51/0.00, -3.45/-0.06, -3.41/-0.12, -3.31/-0.17, -3.25/-0.23, -3.21/-0.28, -3.12/-0.33, -3.09/-0.38, -3.02/-0.42, -2.96/-0.47, -2.92/-0.51, -2.86/-0.56, -2.81/-0.60, -2.76/-0.64, -2.72/-0.68, -2.68/-0.72, -2.66/-0.76, -2.59/-0.79, -2.55/-0.83, -2.49/-0.86, -2.46/-0.90, -2.42/-0.93, -2.36/-0.95, -2.33/-0.99, -2.31/-1.03, -2.28/-1.06, -2.24/-1.09, -2.19/-1.12, -2.16/-1.15, -2.11/-1.17, -2.11/-1.22, -2.04/-1.22, -2.03/-1.27, -1.97/-1.28, -1.95/-1.32, -1.91/-1.34, -1.90/-1.38, -1.86/-1.40, -1.82/-1.42, -1.80/-1.46, -1.76/-1.48, -1.74/-1.51, -1.70/-1.53, -1.67/-1.55, -1.64/-1.58, -1.60/-1.60, -1.58/-1.63, -1.54/-1.65, -1.52/-1.68, -1.49/-1.71, -1.48/-1.76, -1.43/-1.77, -1.40/-1.79, -1.37/-1.81, -1.35/-1.86, -1.31/-1.88, -1.28/-1.89, -1.26/-1.93, -1.22/-1.95, -1.18/-1.96, -1.15/-2.00, -1.11/-2.01, -1.07/-2.01, -1.06/-2.07, -1.02/-2.10, -0.98/-2.11, -0.97/-2.18, -0.92/-2.17, -0.89/-2.20, -0.86/-2.23, -0.82/-2.27, -0.79/-2.30, -0.75/-2.30, -0.72/-2.37, -0.68/-2.36, -0.65/-2.42, -0.61/-2.44, -0.57/-2.48, -0.53/-2.50, -0.49/-2.52, -0.45/-2.57, -0.41/-2.62, -0.37/-2.64, -0.33/-2.70, -0.29/-2.72, -0.24/-2.76, -0.20/-2.79, -0.15/-2.81, -0.10/-2.85, -0.05/-2.91, -0.00/-2.93, 0.05/-2.99, 0.11/-3.01, 0.16/-3.09, 0.22/-3.13, 0.28/-3.18, 0.34/-3.25, 0.40/-3.28, 0.47/-3.37, 0.54/-3.38, 0.61/-3.45, 0.68/-3.50, 0.76/-3.57, 0.84/-3.64, 0.93/-3.72, 1.02/-3.81, 1.11/-3.87, 1.21/-3.96, 2.31/-4.00, 2.34/-3.90, 2.42/-3.87, 2.46/-3.79, 2.52/-3.74, 2.57/-3.68, 2.63/-3.63, 2.69/-3.57, 2.72/-3.48, 2.77/-3.42, 2.84/-3.38, 2.87/-3.30, 2.92/-3.25, 2.98/-3.20, 3.01/-3.12, 3.08/-3.08, 3.13/-3.03, 3.54/-3.30, 4.09/-2.46, 3.71/-2.14, 3.42/-1.89, 3.15/-1.67, 2.92/-1.49, 2.73/-1.33, 2.60/-1.21, 2.46/-1.09, 2.31/-0.98, 2.21/-0.89, 2.12/-0.81, 2.03/-0.74, 1.94/-0.67, 1.87/-0.61, 1.79/-0.55, 1.77/-0.51, 1.69/-0.45, 1.64/-0.41, 1.59/-0.37, 1.54/-0.33, 1.49/-0.29, 1.52/-0.27, 1.52/-0.24, 1.57/-0.22, 1.56/-0.19, 1.62/-0.17, 1.63/-0.14, 1.63/-0.11, 1.68/-0.09, 1.70/-0.06, 1.72/-0.03} \fill[black!25] ({\x - (-4.5) + 10.0},{\y - (-4.0)}) circle (0.9pt);
  \foreach \x/\y in {4.36/3.21, 4.39/3.34, 3.10/4.93, 2.77/4.52, 2.69/4.56, 2.62/4.62, 2.55/4.68, 2.49/4.74, 2.40/4.76, 2.34/4.84, 2.27/4.90, 2.20/4.99, 1.13/2.47, 1.06/2.37, 0.99/2.27, 0.94/2.24, 0.88/2.16, 0.83/2.08, 0.78/2.02, 0.74/2.00, 0.70/1.94, 0.66/1.89, 0.62/1.82, 0.59/1.80, 0.56/1.77, 0.52/1.69, 0.49/1.65, 0.47/1.64, 0.44/1.61, 0.41/1.58, 0.39/1.56, 0.37/1.51, 0.35/1.48, 0.33/1.40, 0.30/1.44, 0.28/1.41, 0.26/1.39, 0.25/1.37, 0.23/1.34, 0.21/1.32, 0.19/1.29, 0.18/1.26, 0.16/1.25, 0.15/1.22, 0.13/1.22, 0.12/1.19, 0.11/1.17, 0.09/1.17, 0.07/1.15, 0.06/1.14, 0.05/1.11, 0.03/1.11, 0.03/1.07, 0.02/1.06, 0.01/1.04, -0.01/1.03, -0.02/1.02, -0.03/0.99, -0.05/1.00, -0.07/1.01, -0.11/1.04, -0.15/1.07, -0.19/1.10, -0.25/1.15, -0.29/1.18, -0.36/1.25, -0.43/1.30, -0.51/1.36, -0.60/1.43, -0.70/1.51, -0.83/1.61, -1.00/1.75, -1.17/1.88, -1.39/2.06, -1.71/2.32, -2.08/2.61, -2.65/3.07, -4.40/0.63, -4.28/0.54, -4.18/0.45, -4.06/0.37, -3.99/0.29, -3.87/0.21, -3.78/0.13, -3.71/0.06, -3.62/-0.01, -3.54/-0.07, -3.45/-0.13, -3.39/-0.19, -3.31/-0.25, -3.24/-0.31, -3.18/-0.37, -3.10/-0.42, -3.04/-0.47, -2.96/-0.51, -2.92/-0.57, -2.87/-0.62, -2.82/-0.66, -2.76/-0.71, -2.70/-0.75, -2.64/-0.79, -2.55/-0.82, -2.54/-0.87, -2.49/-0.91, -2.43/-0.94, -2.39/-0.99, -2.31/-1.01, -2.29/-1.06, -2.24/-1.09, -2.22/-1.14, -2.15/-1.16, -2.12/-1.20, -2.09/-1.24, -2.01/-1.25, -2.00/-1.30, -1.95/-1.33, -1.92/-1.36, -1.88/-1.40, -1.83/-1.43, -1.77/-1.44, -1.74/-1.48, -1.73/-1.52, -1.68/-1.54, -1.65/-1.59, -1.59/-1.59, -1.57/-1.63, -1.54/-1.67, -1.51/-1.71, -1.45/-1.72, -1.42/-1.75, -1.37/-1.77, -1.36/-1.82, -1.32/-1.84, -1.28/-1.87, -1.25/-1.90, -1.21/-1.93, -1.17/-1.95, -1.13/-1.98, -1.09/-2.01, -1.06/-2.04, -1.02/-2.07, -0.98/-2.09, -0.96/-2.14, -0.92/-2.17, -0.88/-2.19, -0.84/-2.21, -0.80/-2.25, -0.77/-2.28, -0.73/-2.31, -0.67/-2.30, -0.66/-2.39, -0.61/-2.41, -0.57/-2.44, -0.53/-2.47, -0.49/-2.51, -0.45/-2.52, -0.41/-2.57, -0.37/-2.62, -0.32/-2.63, -0.28/-2.65, -0.23/-2.68, -0.19/-2.75, -0.14/-2.78, -0.10/-2.80, -0.05/-2.84, -0.00/-2.90, 0.05/-2.92, 0.10/-2.98, 0.15/-3.00, 0.21/-3.03, 0.26/-3.10, 0.32/-3.12, 0.37/-3.19, 0.43/-3.21, 0.49/-3.28, 0.55/-3.34, 0.62/-3.35, 0.69/-3.44, 0.75/-3.50, 0.83/-3.55, 0.90/-3.62, 0.98/-3.67, 1.05/-3.73, 1.14/-3.79, 1.22/-3.85, 1.31/-3.93, 1.40/-3.99, 2.37/-3.96, 2.44/-3.92, 2.49/-3.83, 2.55/-3.78, 2.58/-3.69, 2.66/-3.66, 2.70/-3.58, 2.74/-3.50, 2.81/-3.46, 2.84/-3.38, 2.90/-3.33, 2.94/-3.26, 2.99/-3.20, 3.04/-3.15, 3.10/-3.10, 3.13/-3.03, 3.21/-3.01, 3.24/-2.93, 3.60/-3.20, 4.13/-3.62, 3.77/-2.06, 3.27/-1.66, 2.90/-1.36, 2.58/-1.11, 2.39/-0.95, 2.18/-0.80, 2.07/-0.70, 1.92/-0.59, 1.85/-0.52, 1.75/-0.45, 1.66/-0.38, 1.58/-0.32, 1.54/-0.28, 1.52/-0.25, 1.51/-0.22, 1.55/-0.21, 1.57/-0.20, 1.59/-0.18, 1.59/-0.16, 1.59/-0.14, 1.61/-0.12, 1.64/-0.10, 1.64/-0.08, 1.69/-0.07, 1.70/-0.05, 1.72/-0.03, 1.72/-0.00, 1.75/0.02, 1.73/0.04, 1.78/0.06, 1.80/0.08, 1.80/0.11, 1.82/0.13, 1.84/0.16, 1.86/0.18, 1.89/0.21, 1.89/0.23, 1.93/0.26, 1.93/0.29, 1.96/0.32, 2.00/0.35, 2.03/0.38, 2.02/0.41, 2.10/0.45, 2.12/0.49, 2.15/0.52, 2.19/0.57, 2.19/0.60, 2.25/0.65, 2.28/0.69, 2.30/0.73, 2.36/0.79, 2.41/0.84, 2.45/0.89, 2.45/0.93, 2.52/1.01, 2.57/1.07, 2.62/1.13, 2.68/1.21, 2.76/1.29, 4.46/3.17} \fill[ohmamber] ({\x - (-4.5) + 10.0},{\y - (-4.0)}) circle (0.8pt);
  \foreach \x/\y in {4.43/3.11, 4.45/3.24, 3.20/4.86, 2.86/4.46, 2.78/4.50, 2.71/4.55, 2.65/4.61, 2.58/4.68, 2.50/4.70, 2.44/4.78, 2.37/4.84, 2.31/4.93, 2.22/4.96, 1.18/2.43, 1.10/2.34, 1.03/2.24, 0.98/2.21, 0.92/2.13, 0.87/2.05, 0.82/2.00, 0.78/1.98, 0.74/1.91, 0.70/1.87, 0.66/1.80, 0.62/1.78, 0.59/1.75, 0.55/1.67, 0.52/1.63, 0.50/1.62, 0.47/1.60, 0.44/1.56, 0.42/1.54, 0.40/1.49, 0.37/1.46, 0.35/1.39, 0.33/1.43, 0.31/1.39, 0.29/1.37, 0.27/1.35, 0.25/1.33, 0.23/1.31, 0.22/1.28, 0.20/1.25, 0.19/1.23, 0.17/1.21, 0.15/1.21, 0.14/1.18, 0.13/1.16, 0.11/1.16, 0.10/1.14, 0.08/1.13, 0.07/1.10, 0.05/1.10, 0.05/1.07, 0.03/1.05, 0.02/1.03, 0.01/1.02, -0.00/1.01, -0.01/0.99, -0.03/1.00, -0.05/1.00, -0.09/1.03, -0.13/1.07, -0.17/1.09, -0.23/1.15, -0.27/1.18, -0.34/1.24, -0.41/1.30, -0.48/1.36, -0.57/1.44, -0.67/1.52, -0.80/1.62, -0.96/1.76, -1.13/1.90, -1.35/2.08, -1.66/2.35, -2.02/2.64, -2.59/3.12, -4.39/0.72, -4.27/0.62, -4.18/0.53, -4.06/0.44, -3.98/0.36, -3.87/0.28, -3.78/0.21, -3.71/0.13, -3.63/0.06, -3.54/-0.00, -3.46/-0.07, -3.39/-0.13, -3.32/-0.19, -3.25/-0.25, -3.19/-0.31, -3.12/-0.36, -3.05/-0.41, -2.98/-0.46, -2.94/-0.51, -2.89/-0.56, -2.84/-0.61, -2.78/-0.66, -2.72/-0.70, -2.66/-0.74, -2.57/-0.77, -2.56/-0.83, -2.51/-0.87, -2.45/-0.90, -2.41/-0.94, -2.33/-0.97, -2.32/-1.02, -2.27/-1.05, -2.25/-1.10, -2.18/-1.12, -2.15/-1.16, -2.12/-1.20, -2.04/-1.22, -2.03/-1.27, -1.98/-1.29, -1.95/-1.33, -1.92/-1.37, -1.87/-1.39, -1.81/-1.41, -1.78/-1.45, -1.76/-1.50, -1.71/-1.51, -1.69/-1.56, -1.63/-1.57, -1.60/-1.61, -1.58/-1.65, -1.55/-1.69, -1.49/-1.69, -1.46/-1.72, -1.41/-1.75, -1.40/-1.80, -1.36/-1.82, -1.32/-1.85, -1.29/-1.88, -1.25/-1.91, -1.21/-1.93, -1.18/-1.96, -1.14/-1.99, -1.11/-2.03, -1.07/-2.05, -1.03/-2.07, -1.01/-2.12, -0.97/-2.15, -0.93/-2.18, -0.89/-2.20, -0.85/-2.24, -0.82/-2.27, -0.78/-2.30, -0.72/-2.29, -0.71/-2.39, -0.67/-2.40, -0.63/-2.44, -0.59/-2.47, -0.55/-2.51, -0.50/-2.52, -0.47/-2.57, -0.43/-2.62, -0.38/-2.63, -0.34/-2.66, -0.29/-2.68, -0.25/-2.75, -0.21/-2.78, -0.16/-2.81, -0.11/-2.85, -0.07/-2.91, -0.02/-2.93, 0.03/-2.99, 0.08/-3.01, 0.14/-3.04, 0.19/-3.11, 0.25/-3.13, 0.30/-3.20, 0.36/-3.23, 0.42/-3.30, 0.48/-3.36, 0.54/-3.37, 0.61/-3.47, 0.68/-3.52, 0.75/-3.57, 0.82/-3.65, 0.89/-3.70, 0.97/-3.76, 1.05/-3.82, 1.13/-3.88, 1.22/-3.97, 2.35/-3.98, 2.40/-3.89, 2.46/-3.84, 2.50/-3.75, 2.58/-3.72, 2.62/-3.65, 2.66/-3.57, 2.73/-3.53, 2.77/-3.45, 2.82/-3.40, 2.86/-3.32, 2.92/-3.28, 2.97/-3.22, 3.03/-3.17, 3.06/-3.10, 3.14/-3.08, 3.17/-3.00, 3.53/-3.29, 4.05/-3.72, 3.72/-2.14, 3.23/-1.73, 2.86/-1.43, 2.55/-1.18, 2.36/-1.01, 2.16/-0.85, 2.05/-0.75, 1.91/-0.64, 1.83/-0.57, 1.74/-0.49, 1.65/-0.42, 1.57/-0.36, 1.53/-0.32, 1.51/-0.29, 1.50/-0.26, 1.54/-0.25, 1.56/-0.24, 1.58/-0.22, 1.58/-0.20, 1.58/-0.18, 1.60/-0.16, 1.63/-0.15, 1.64/-0.13, 1.69/-0.11, 1.69/-0.09, 1.71/-0.07, 1.72/-0.05, 1.75/-0.03, 1.72/-0.00, 1.78/0.02, 1.80/0.04, 1.80/0.06, 1.81/0.09, 1.84/0.11, 1.86/0.13, 1.90/0.16, 1.89/0.19, 1.93/0.21, 1.93/0.24, 1.97/0.27, 2.00/0.30, 2.04/0.33, 2.02/0.36, 2.10/0.40, 2.13/0.44, 2.15/0.47, 2.20/0.51, 2.20/0.54, 2.26/0.59, 2.29/0.63, 2.31/0.68, 2.37/0.73, 2.43/0.78, 2.46/0.83, 2.46/0.87, 2.54/0.94, 2.59/1.00, 2.64/1.07, 2.70/1.14, 2.78/1.23} \fill[ohmred] ({\x - (-4.5) + 10.0},{\y - (-4.0)}) circle (0.8pt);
  \fill[white] (14.50,4.00) circle (2pt);
  \node[below right, white, font=\scriptsize] at (10.00,9.00) {task 2};
  \end{scope}
\end{tikzpicture}
}
  \caption{One match 20 scans apart in a corner, in the frame of the robot at the old scan (white dot). Grey
    the model, amber the new scan where ICP starts, red where ICP put it. Left, task 1, from the identity:
    0.5\,m off. Right, task 2, from the odometry: 1\,cm off.}
  \label{fig:match}
\end{figure}

\expected[\thema{every:=20}; ]{every20}

\begin{checkpoint}
Explain with \cref{tab:every}: too close and the small errors of many matches add up; too far and ICP starts
outside the valley of the right minimum. The script calls this convergence to a \emph{local} minimum.
\end{checkpoint}

\newpage
\section{Task 2: ICP with odometry\minutes{45}}
\label{sec:task2}
\running{}{-p task:=2 -p every:=2 \textsf{, then 5, 10 and 20}}

ICP alone starts every match at the identity, as if the robot had not moved. The odometry knows better: it
measured how far the robot went and how much it turned. Start ICP there. The odometry gives two poses in its
own frame, \(\matsym T_\text{before}\) at the old scan and \(\matsym T_\text{now}\). The motion in between,
in the frame of the robot at the old scan:
\begin{equation*}
  \matsym T_0 = \matsym T_\text{before}^{-1}\, \matsym T_\text{now} .
\end{equation*}
\thema{icp(M, D, T\_init)} moves \(D\) by \(\matsym T_0\) first and starts the product with it. The rest of
ICP stays as it is.

\begin{todo}{2.1}
\datei{scan\_odometry.py}, \thema{odom\_guess(odom\_prev, odom\_now)}: \(\matsym T_0\) from two
\thema{nav\_msgs/msg/Odometry}. \thema{se2} and \thema{yaw} are given.
\end{todo}

Start the node with \thema{task:=2} and try the values of \cref{tab:odom}, one lap each.

\begin{table}[ht]\centering
\renewcommand{\arraystretch}{1.4}
\caption{One lap, ICP with the odometry as the start. Fill in. ms: per match, terminal 2.}
\label{tab:odom}
\begin{tabular}{@{}p{2.6cm}p{2.0cm}p{2.0cm}p{2.0cm}p{2.0cm}p{2.0cm}@{}}
\toprule
\thema{every}, \thema{d\_max} & driven in m & turned in \textdegree & error in m & error in \textdegree & ms \\
\midrule
2, 0.3 & & & & & \\
5, 0.3 & & & & & \\
10, 0.3 & & & & & \\
20, 0.3 & & & & & \\
10, 0 & & & & & \\
\bottomrule
\end{tabular}
\end{table}

\begin{itemize}
  \item \thema{every} = 20 failed in task 1. Why does it work now (\cref{fig:match}, right)? Why does ICP need
        fewer iterations?
  \item The odometry drives 3\,\% too far. ICP starts from it. Does the ICP path come out 3\,\% too long?
  \item \thema{every} = 2 is still too short, even from the odometry. Why does a good start not help here?
  \item Without \thema{d\_max} (last row) ICP is about as bad in heading as the odometry it started from,
        or worse. Why?
  \item Which sensor knows the distance well, which the heading? Which \thema{every} do you keep?
\end{itemize}

\expected[\thema{every:=10}; ]{odom}

\begin{checkpoint}
\thema{pytest -k task2} passes. After one lap ICP with odometry is more than five times better than the
odometry, in position and in heading. Show the plots of \thema{plot\_run.py} for task 1 and task 2.
\end{checkpoint}

\newpage
\section{Extra task: the kd-tree\minutes{15}}
\label{sec:extra}
\running{}{-p task:=1 -p kdtree:=true}

Step 1 is the expensive one. Brute force compares every point of \(D\) with every point of \(M\):
\(\mathcal O(N_d N_m)\) per iteration, here about \(300 \times 300\). A kd-tree of \(M\) is built once per match,
\(\mathcal O(N_m \log N_m)\); then each query costs \(\mathcal O(\log N_m)\) (script, section 4.1.3).
\thema{scipy.spatial.cKDTree} is such a tree. With \thema{-p kdtree:=true} the node builds it and calls your
function instead of \thema{closest\_points}.

\begin{todo}{Z.1}
\datei{icp.py}, \thema{closest\_points\_kdtree(D, tree)}: the same \thema{(index, distance)} as
\thema{closest\_points}, with one call of \thema{tree.query(D)}.
\end{todo}

\begin{itemize}
  \item Compare the ms per match in terminal 2 with \cref{tab:every} (\thema{every} = 10). How many times
        faster? Is the pose the same?
  \item Brute force needs \(N_d N_m\) distances per iteration. How many is that for two scans of 300 points?
        And for two 3D scans of 100\,000 points each?
\end{itemize}

\newpage
\section{Extra task: every scan\hfill{\normalsize\mdseries\color{ohmgrey}if time is left}}
\label{sec:extra2}
\running{}{-p task:=1 -p kdtree:=true -p keyframe:=true}

In \cref{sec:task1d} the small steps failed: \thema{every} = 2 chains five times as many matches as 10, and the
small error of each match adds up. So do not chain every scan onto the one before. Keep one scan as the model,
the \emph{keyframe}, and match every new scan onto it, 20 times a second. Only when the robot is
\thema{key\_dist} away from the keyframe or turned by \thema{key\_angle}, the newest scan becomes the new
keyframe. A pose 20 times a second, but a new link in the chain only every half metre:
\begin{equation*}
  \matsym T_\text{robot} = \matsym T_\text{keyframe} \, \matsym T , \qquad
  \matsym T \text{ maps the newest scan onto the keyframe.}
\end{equation*}
In task 1 ICP starts at the last match onto the same keyframe: the robot moved only 2.5\,cm since. The node
does the bookkeeping. It needs the kd-tree of \cref{sec:extra}: brute force takes longer than the 50\,ms
between two scans.

\begin{todo}{Z.2}
\datei{scan\_odometry.py}, \thema{new\_keyframe(T, key\_dist, key\_angle)}: \thema{True} when the robot is
farther than \thema{key\_dist} from the keyframe or turned by more than \thema{key\_angle} (radians), in
either direction.
\end{todo}

The default is 0.5\,m and 20\,\textdegree. Run one lap each with the values of \cref{tab:keyframe}
(\thema{-p key\_dist:=1.0 -p key\_angle:=30.0}).

\begin{table}[ht]\centering
\renewcommand{\arraystretch}{1.4}
\caption{One lap, every scan onto the keyframe, ICP alone. Fill in. Keyframes: terminal 2.}
\label{tab:keyframe}
\begin{tabular}{@{}p{2.6cm}p{2.0cm}p{2.0cm}p{2.0cm}p{2.0cm}p{2.0cm}@{}}
\toprule
\thema{key\_dist}, \thema{key\_angle} & driven in m & turned in \textdegree & error in m & error in \textdegree & keyframes \\
\midrule
0.5, 20 & & & & & \\
1.0, 30 & & & & & \\
3.0, 90 & & & & & \\
5.0, 180 & & & & & \\
\bottomrule
\end{tabular}
\end{table}

\begin{itemize}
  \item Compare with \cref{tab:every}. Why is the path no longer too short?
  \item In \cref{sec:task1d} half a metre between two scans was too much for ICP alone. Now 1\,m works. What is
        different? (Hint: \cref{fig:match}.)
  \item Why does 3\,m still work, and 5\,m and 180\,\textdegree{} not? What must the keyframe and the newest
        scan have in common?
\end{itemize}

\expected[\thema{keyframe:=true}; ]{keyframe}

\section{Outlook: a map}
\label{sec:outlook}
Your ICP chains scan to scan, or keyframe to keyframe. Each link adds a small error, and nothing ever takes it
back: it is odometry, only a much better one. Match the scan against a \emph{map} instead, built from all scans so far, and the
error stops growing while the robot stays in known places. When it comes back after a long loop and the scan
does not fit the map, the robot has found its own drift: that is loop closure, and with it SLAM (script,
sections 4.2 to 4.4).

\end{document}
